\section{Software Tools}
\label{ch:SoftwareTools}

In this section, we describe the tool and script developed, along with the approaches used,
to create the two datasets described in \cref{sec:datasets}.

\subsection{Tool for Analyzing NPM Packages}
\label{sec:ToolForAnalyzingNpmPackages}

The \texttt{Analyzer-npm-package-releases} tool computes
the metrics of \cref{ch:metricsAnalysis} from a set of \ac{NPM} package versions.

Using this tool, it is possible to construct 
the Goodware dataset (\cref{sec:goodwareDataset}) and the Malware dataset (\cref{sec:malwareDataset}).

Implemented in Python, it is available on GitHub\footnote{\url{https://github.com/Frazzerz/Analyzer-npm-package-releases}},
and can be invoked from the command line in the following way:

\begin{verbatim}
python3 main.py --json packages.json [--workers N] [--local]
\end{verbatim}

The tool requires as input a JSON file containing a list of \ac{NPM} packages names you want to analyze.

Instead, the \texttt{--local} flag allows local tarballs of \ac{NPM} package versions to be include in the analysis.

The tool processes one package and one version at a time, from the oldest version to the most recent one.
Instead, the analysis and computation of metrics for individual files within in a version are performed in parallel, using multiprocessing.
The \texttt{--workers} flag allows for specifying the number of processes to be used.

The tool's workflow can be divided into three phases:
\begin{enumerate}
    \item \textbf{Version Retrieval}:

First, the NPMClient component retrieves the package metadata and the list of available versions
by querying the official \ac{NPM} registry.

From this list, only versions parsable via node-semver \cite{semver} are kept, 
to make them chronologically sortable.

The tool selects the 20 most recent valid versions from the list and downloads 
the related tarballs for the Goodware dataset (\cref{sec:goodwareDataset}).

Instead, the tool selects and downloads the 19 most recent valid tarballs versions
for the Malware dataset (\cref{sec:malwareDataset}).
In this case, the twentieth version is the tarball related to the malicious version,
retrieved with the developed script (\cref{sec:DownloadMaliciousTarballs}) 
and included using the \texttt{--local} flag.

    \item \textbf{Metrics Computation}:

In this phase, the tool first computes the metrics on the individual files of the package version's tarball.

The computation of each metric is different, and each one is described in 
depth in its own section in \cref{ch:metricsAnalysis}.

Metrics based on the textual content of the files,
such as Unique Hexadecimal values (\cref{sec:analysisHexUnique}),
are computed only on plain text files, such as JavaScript, TypeScript, JSON and Markdown files,
ignoring binary files such as zip, png or dll.

Furthermore, for these metrics, the tool removes comments 
from JavaScript and TypeScript files using a parser from the tree-sitter library \cite{treesitter}.

Subsequently, these metrics computed on individual files are aggregated to form metrics at the package version level.

For example, to compute the package size metric (\cref{sec:analysisPackageSize}),
the tool sums the sizes of all files contained in the tarball.

Aggregation strategies vary based on the type of metric. 
For instance, counts and weights are summed, 
the maximum value among the longest lines of the files is taken,
and strings are collected into lists.

    \item \textbf{CSV Reporter}:

Finally, for each package, the results are saved incrementally in two CSV files.
\begin{itemize}
    \item \texttt{file\_metrics.csv} contains one row for every file of each package version. 
    The columns represent the metrics computed on single file.
    \item \texttt{aggregate\_metrics\_by\_single\_version.csv} contains one row for every package version, 
    and the columns represent the metrics computed at the package version level.
\end{itemize}

The tool saves the results incrementally at the completion of each version's analysis,
ensuring that partial results are preserved even the event of an error.

In addition to the CSV files, 
the tool produces a log.txt log file, which is also updated incrementally
and serves as a copy of the terminal output.

Inside the file, there is general information such as the number of packages to analyze and the number of workers that will be used.
Furthermore, for each package, the file also reports: the number of versions found in the registry,
the status of the tarball downloads, the progress of the analysis of the selected versions
and the total time for the analysis of the entire package.

\end{enumerate}